%%%%%%%%%%%%%%%%%%%%%%%%%%%%%%%%%%%%%%%%%%%%%%%%%%%%%
%			HLAVIČKA								%
%%%%%%%%%%%%%%%%%%%%%%%%%%%%%%%%%%%%%%%%%%%%%%%%%%%%%
\documentclass[openany,12pt]{memoir}

\usepackage[utf8]{inputenc} 
\usepackage[czech]{babel}
\usepackage[T1]{fontenc}
\usepackage[top=1.5cm, bottom=2cm, left=2cm, right=2cm]{geometry}  % --> NASTAVENÍ OKRAJŮ
\usepackage{fancyhdr}
\usepackage{graphicx}
\usepackage{lmodern}
\usepackage{xwatermark}
\usepackage{xcolor}
\usepackage{changepage}
\usepackage{newunicodechar}



%%%%%% Package na zpěvník
\usepackage[full]{leadsheets}%http://mirrors.nic.cz/tex-archive/macros/latex/contrib/leadsheets/leadsheets_en.pdf   --> dokumentace	
\definesongtitletemplate{empty}{} 
\setchords{
format = \bfseries,   %tučné akordy
minor = {mi},% 
input-notation = {german},%
output-notation = {german}%
}
\definesongtitletemplate{empty}{} 

\newlength{\drop}
\newwatermark[allpages,color=red!50,angle=0,scale=2, xpos=0,ypos=0]{\includegraphics[width=5cm]{obr/pozadi.jpg}} %--> dvojka na pozadí


%%%% Vlastní příkazy
\newcounter{Slokočet}   %Automatické číslování slok
\newcommand{\mezera}{\vspace*{0.5cm}}   %Horizontální odsazení slok
\newcommand{\stred}{5.2cm}   %%% Na zarovnání slok doprostřed, pozn. automatičtější zarovnávání na střed nejde
\newcommand{\refren}{\mezera \noindent \textbf{R:} } %refrén
\newcommand{\sloka}{\mezera \noindent \addtocounter{Slokočet}{1} \arabic{Slokočet}. } 	%sloka, která se automaticky čísluje
\newcommand{\ssloka}{\mezera \noindent}  % vlastní číslo sloky
\newcommand{\carka}{,\:}
\newcommand{\m}[1]{\color{white}{#1}}  %Pro akordy

%%%%%%%%%%%%%%%%%%%%%%%%%%%%%%%%%%%%%%%%%%%%%%%%%%%%%%
%			NÁVOD									 %
%%%%%%%%%%%%%%%%%%%%%%%%%%%%%%%%%%%%%%%%%%%%%%%%%%%%%%
%1. Věci v hlavičce IGNOROVAT
%2. Píseň psát do prostoru mezi \begin{song} a \end{song}
%3. další řádek se značí dvěma odsazeníma (= dvakrát stisknout enter)
%4. \refren vždy na začátku refrenu a \sloka na začátek sloky (automaticky se čísluje)
%  \ = alt gr + q ; [/] = alt gr f/g ; {/} = alt gr + b/n; ^ = alt gr + 3 + mezera
%Cokoliv napíšete do ^{  } se bude brát jako akord
%když se toto bude dotýkat nějakého slova (nebude mezi tím a slovem mezera)
%tak se akordy zjeví nad slovem, ale pište to před slovo
%Když se to nedotýká slova, tak akord lítá ve vzduchu a vytiskne se větší mezera
%První možnost je asi preferovanější
%5. Akordy stačí psát jen do první sloky, když se nezmění -- kytaristi to zvládnou



%<++>
%\usepackage{subfiles}
%</++>

\begin{document}

\pagestyle{empty}
%\addtocontents{toc}{\protect\thispagestyle{empty}}
%\tableofcontents \thispagestyle{empty}\newpage
\newgeometry{top=1.5cm, bottom = 0cm, left = 2cm, right = 2cm}



\begin{song}{title=\centering Otevřená zlomenina srdečního svalu \\\normalsize Wanastowi Vjecy \vspace*{-0.3cm}}  %% sem se napíše jméno songu a autor
\moveright 4.5cm \vbox{      %Varianta č. 1  ---> Jeden sloupec zarovnaný na střed	

\sloka 
	Jsem jako ^{G}vítr kterej zfoukne pírko ze tvejch dlaní,

	^{D}špínu jedný noci, jako hygiena ranní,

	^{Ami}velká voda slzí, který spláchnou noční hříchy, 

	^{C}tamburíny zvoní k ^{Es}operaci míchy.
	
	Narovnám ti páteř, poškrábu ti záda, 
	
	pofoukám ti srdce, zrada kamaráda

	lásky účel světí prostředky a smetí, 
	
	špínu jedný noci, který neuletíš.

	Otevřená zlomenina srdečního svalu, 
	
	trápení a kocovina, vůně tvýho žalu.
	
	Vypustíš svou hříšnou duši do slanýho moře, 
	
	neumírej děvče moje chti ti říct že:


\refren
	Ohořelou ^{G}károu chtěl bych dojet ^{D}ke hvězdám,

	který svítily z tvejch ^{Ami}očí dřív než ^{C}červotoči 

	se do tvýho srdce ^{F}daj, hm ^{D}hm.
	
	V ohořelym autě už dva měsíce nedejchám,
	
	sám se svojí vinou, už nikdy nechci jinou, 
	
	už asi nedoufám.

\sloka
	Pláčem solíš otevřený rány co se hojí, 
	
	tvoji krev i tělo příjímám pod obojí.

	Stejný lidi se soumrakem mají stejný stíny, 
	
	zmizelas jak před přízrakem, padám do hlubiny.
	
	Tvůj pramínek vlasů zaliju včelím voskem, 
	
	koukám na tu krásu a nechápu to mozkem.
	
	To co jsi mi dala já nikomu už nedám, 
	
	tak mi řekni má opičko proč tě marně hledám.

	Otevřená zlomenina srdečního svalu, 
	
	trápení a kocovina, vůně tvýho žalu.
	
	Vypustíš svou hříšnou duši do slanýno moře, 

	neumírej děvče moje chci ti říct, že:

\refren




}
\setcounter{Slokočet}{0}
\end{song}

\newpage


\end{document}
